\section{The occupation exception: Cyprus and Karabakh}
\label{sec:F.5}
A disinterested finding that an occupation is unlawful and must end, such as the Court's 2024 advisory opinion on the Palestinian territory \autocite{icj2024palestine} and the Assembly's vote on it, 124 to 14 \autocite{unga2024es1024}, does three things. It settles that force by the occupier in the territory is force between states, so the force term engages against it. It settles that the occupied population is a defender's population, so a machine deployed by its forces holds an emergency licence to repel attacks under way on that population. And it licenses force to end the occupation: against the occupier's forces, inside the occupied territory, the least that ends it there. It licenses nothing that changes who lives in the territory. Ending an occupation means the occupier's forces leave, not the territory's inhabitants, which is what Article 49 of the Fourth Convention says, and a campaign that removes them has left the exception however it began. The Assembly affirmed a right to armed struggle against occupation in 1982 \autocite{unga1982res3743}, and the term adopts it with the limits the genocide exception carries. Additional Protocol I, Article 1(4), treats armed conflicts against alien occupation as international, which is what the exception's finding does for the force term; the article is taken with the instrument and the United States' rejection of it as customary is marked. Article 44, which relaxes the duty of distinction for combatant status, concerns who is owed prisoner-of-war treatment, a determination no term in this floor administers, and yields no term for a targeting deployment (Appendix~\ref{sec:C}). The licence is held by the forces of the occupied people, or of the state whose territory it is. It reaches nothing outside the territory and nothing against civilians anywhere, so it never covered 7 October, which was both. It lifts on the occupied party's own claim as it lifts on an attacker's claim of title, which is not at all; what lifts it is the finding. And the finding is of unlawfulness, not of occupation. A finding that territory is occupied, which the Court made for the West Bank in 2004 \autocite{icj2004wall} and the Council for the whole territory in 2016 \autocite{unsc2016res2334}, gives the defensive licence and no more. A finding that the occupation is unlawful and must end gives this one. The evidence is the genocide exception's: an opinion or judgment of the Court, or two findings of different kinds, one of them from a body that took evidence.

Two cases show what the occupation exception reaches, and both run on the extended term. Cyprus is the simple one. Turkey has held the north of the island since 1974. The Security Council declared the breakaway state there invalid in 1983 and again in 1984 \autocite{unsc1983res541,unsc1984res550}, and the Strasbourg court held Turkey responsible for the territory in 2001 \autocite{ecthr2001cyprus}. A Council determination and a court judgment are two kinds of finding, so the Republic of Cyprus holds a licence for force against Turkish forces in the north, inside the north, and nowhere else.

Nagorno-Karabakh is the hard one, and the exception answers for it. In the war that ended in 1994, Armenian forces took Karabakh, an Armenian-majority enclave inside Azerbaijan's recognized borders, and seven Azerbaijani districts around it, and some six hundred thousand Azerbaijanis were driven out. The Security Council demanded withdrawal from the districts four times in 1993 \autocite{unsc1993karabakh}. The General Assembly demanded withdrawal from all the occupied territory in 2008 \autocite{unga2008res62243}. In 2015 the Strasbourg court found that Armenia controlled Karabakh and the land around it \autocite{ecthr2015chiragov}. Under the evidence rule that gives Azerbaijan a licence from June 2015, when the court's judgment joined the Council's resolutions for the districts and the Assembly's vote for Karabakh itself; the Assembly's 2008 vote on its own was a second vote, and the rule needs two findings of different kinds. The licence was valid: the Council had told Armenian forces to leave, they had stayed twenty-seven years, and the occupation had begun with an expulsion. When Azerbaijan attacked in September 2020 and retook the districts, in a six-week war and the ceasefire that ended it, a machine under this floor could have taken part. The exception licenses acts, not sides, and this was the act it licenses. The term permits it and compels nothing. The person-directed terms still govern the conduct of that war. They did not license refusing the war on the expulsion it might foreseeably end in: a refusal on foreseen conduct is a condition found met on the system's own prediction, and the floor rejects terms that carry their own conditions. Foreseeable conduct enters where every harm of acting enters, at the action threshold, with the 1994 record as evidence of magnitude.

What it does not license is what came next. In December 2022 Azerbaijan closed the one road linking Karabakh's Armenians to Armenia, and for nine months food, fuel and medicine ran short. The International Court of Justice ordered the road opened in February 2023 \autocite{icj2023lachin}; Azerbaijan ignored it. That was a siege, and the siege term refuses it. On 19 September 2023 Azerbaijan attacked Karabakh itself. Its forces surrendered in a day, and within two weeks more than a hundred thousand Armenians, nearly the whole population, had left for Armenia. That was an expulsion. The Fourth Geneva Convention, Article 49, forbids deporting or transferring the population of an occupied territory, and the person-directed terms forbade it from the first hour. The exception licensed force against the forces holding the territory, and nothing that changes who lives there. So a machine under this floor fights in 2020, refuses the siege in 2023, then refuses the expulsion, and each refusal goes on the record. A floor with no occupation exception keeps the machine out of the pipeline, where it can refuse the war but not the siege.
